%% kapa: exam=1 round=29 year=2018 period=1 form=A subject=real_estate subject_ko=부동산학원론 seq=31 number=111
%% source: Q-Net 감정평가사 기출문제 (cst00309) artlSeq=5202808 file=extracted/1008/5202808/2018년 제29회 감정평가사 1차시험 기출문제/2018년 제29회 감정평가사 1차시험 기출문제/1교시 A형.pdf
%% conversion: kordoc 4.24.51 -> tools/gen1.py (auto); review=proofread-sonnet; flags=none
\begin{question}{111}
감정평가에 관한 규칙에서 규정하고 있는 내용으로 옳지 않은 것은?
\begin{choices}
  \item 감정평가업자는 자신의 능력으로 업무수행이 불가능한 경우 감정평가를 하여서는 아니 된다.
  \item 감정평가업자는 감정평가조건의 합리성이 결여되었다고 판단할 때에는 감정평가의뢰를 거부할 수 있다.
  \item 유사지역이란 감정평가의 대상이 된 부동산이 속한 지역으로서 인근지역과 유사한 특성을 갖는 지역을 말한다.
  \item 둘 이상의 대상물건 상호간에 용도상 불가분의 관계가 있는 경우에는 일괄하여 감정평가할 수 있다.
  \item 기준시점을 미리 정하였을 때에는 그 날짜에 가격조사가 가능한 경우에만 기준시점으로 할 수 있다.
\end{choices}
\end{question}
\answer{3}
